% preamble-exam-zh.tex — 共享 preamble
% 用于所有 exam-zh 模板
%
% 样式策略：
%   屏幕 → 语义彩色（蓝=关键想法，红=易错，绿=路线，灰=提示/教师备注）
%   打印 → 自动降级为黑白（通过 print mode 或 monochrome 选项）

\documentclass{exam-zh}

% ---------- 基础配置 ----------
\examsetup{
  page/size = a4paper,
  page/show-head = false,
  page/show-foot = false,
  sealline/show = false,
  font = times,
  math-font = xits,
}

\usepackage{tcolorbox}
\usepackage{needspace}
\tcbuselibrary{skins,breakable}

% ---------- 打印/屏幕颜色切换 ----------
% 用法：\documentclass[monochrome]{exam-zh} 即可切换为纯黑白
% 注意：不要使用 \if@xxx 放在 makeatother 外部；这里改成普通条件名。
\newif\ifeduprintcolor
\eduprintcolortrue

\makeatletter
\@ifclasswith{exam-zh}{monochrome}{\eduprintcolorfalse}{}
\makeatother

% ---------- 语义颜色定义 ----------
% 屏幕：有色彩区分；打印：降级为灰度
\ifeduprintcolor
  % --- 屏幕彩色模式 ---
  \definecolor{edu-blue}{RGB}{47, 82, 122}       % 关键想法
  \definecolor{edu-blue-bg}{RGB}{243, 246, 252}
  \definecolor{edu-red}{RGB}{180, 40, 40}         % 易错提醒
  \definecolor{edu-red-bg}{RGB}{255, 245, 240}
  \definecolor{edu-green}{RGB}{34, 100, 60}       % 路线图
  \definecolor{edu-green-bg}{RGB}{242, 250, 245}
  \definecolor{edu-orange}{RGB}{160, 90, 0}       % 训练目标
  \definecolor{edu-orange-bg}{RGB}{255, 248, 235}
  \definecolor{edu-gray}{RGB}{100, 100, 100}      % 提示/教师备注
  \definecolor{edu-gray-bg}{RGB}{248, 248, 248}
  \definecolor{edu-purple}{RGB}{90, 50, 120}      % 思考题
  \definecolor{edu-purple-bg}{RGB}{248, 244, 252}
\else
  % --- 打印黑白模式 ---
  \definecolor{edu-blue}{gray}{0.15}
  \definecolor{edu-blue-bg}{gray}{0.96}
  \definecolor{edu-red}{gray}{0.25}
  \definecolor{edu-red-bg}{gray}{0.97}
  \definecolor{edu-green}{gray}{0.20}
  \definecolor{edu-green-bg}{gray}{0.97}
  \definecolor{edu-orange}{gray}{0.25}
  \definecolor{edu-orange-bg}{gray}{0.98}
  \definecolor{edu-gray}{gray}{0.40}
  \definecolor{edu-gray-bg}{gray}{0.97}
  \definecolor{edu-purple}{gray}{0.30}
  \definecolor{edu-purple-bg}{gray}{0.98}
\fi

% ---------- 自定义教学环境 ----------
% 共性：enhanced + breakable（跨页不断裂）

% 原题卡片 — 强边框，突出显示
\newtcolorbox{problemcard}{
  enhanced, breakable,
  colback=edu-blue-bg,
  colframe=edu-blue,
  boxrule=1.2pt,
  arc=0pt,
  left=3mm, right=3mm, top=3mm, bottom=3mm,
}

% 解题路线图 — 左边线 + 浅底
\newtcolorbox{routebox}{
  enhanced, breakable,
  colback=edu-green-bg,
  colframe=edu-green!30,
  boxrule=0pt,
  borderline west={2.5pt}{0pt}{edu-green},
  left=3mm, right=3mm, top=3mm, bottom=3mm,
}

% 关键想法 — 蓝色左边线
\newtcolorbox{keyideabox}{
  enhanced, breakable,
  colback=edu-blue-bg,
  colframe=edu-blue!30,
  boxrule=0pt,
  borderline west={2.5pt}{0pt}{edu-blue},
  left=3mm, right=3mm, top=3mm, bottom=3mm,
}

% 易错提醒 — 红色左边线
\newtcolorbox{mistakebox}{
  enhanced, breakable,
  colback=edu-red-bg,
  colframe=edu-red!20,
  boxrule=0pt,
  borderline west={3pt}{0pt}{edu-red},
  left=3mm, right=3mm, top=2.5mm, bottom=2.5mm,
}

% 提示 — 灰色虚线左边线
\newtcolorbox{hintbox}{
  enhanced, breakable,
  colback=edu-gray-bg,
  colframe=edu-gray!20,
  boxrule=0pt,
  borderline west={2pt}{0pt}{edu-gray, dashed},
  left=3mm, right=3mm, top=2mm, bottom=2mm,
}

% 思考题 — 紫色虚线左边线
\newtcolorbox{thinkbox}{
  enhanced, breakable,
  colback=edu-purple-bg,
  colframe=edu-purple!20,
  boxrule=0pt,
  borderline west={2pt}{0pt}{edu-purple, dashed},
  left=2.5mm, right=2.5mm, top=2.5mm, bottom=2.5mm,
}

% 教师备注 — 灰色左边线，小字
\newtcolorbox{teachernote}{
  enhanced, breakable,
  colback=edu-gray-bg,
  colframe=edu-gray!15,
  boxrule=0pt,
  borderline west={2pt}{0pt}{edu-gray!60},
  left=2.5mm, right=2.5mm, top=2.5mm, bottom=2.5mm,
  fontupper=\small,
  colupper=black!60,
}

% 训练目标 — 橙色左边线，小字
\newtcolorbox{traininggoal}{
  enhanced, breakable,
  colback=edu-orange-bg,
  colframe=edu-orange!15,
  boxrule=0pt,
  borderline west={1.5pt}{0pt}{edu-orange},
  left=2mm, right=2mm, top=1mm, bottom=1mm,
  fontupper=\small,
}

% 答题区 — 无边框，纯留白
\newcommand{\answerarea}[1][40mm]{%
  \vspace{2mm}%
  \begin{tcolorbox}[
    enhanced, breakable,
    colback=white,
    colframe=white,
    boxrule=0pt,
    height=#1,
    valign=top,
    bottom=0mm,
    after skip=0mm,
  ]
  \end{tcolorbox}%
}

% ---------- 字体设置 ----------
% exam-zh (ctexbook) already sets CJK fonts via ctex.
% Only override if specific fonts are needed.
% macOS: Songti SC, STHeiti, STFangsong
% Linux: SimSun, SimHei, FangSong

% ---------- 页面设置 ----------
% exam-zh (ctexbook) already loads geometry.
% Override margins if needed:
% \geometry{top=16mm, bottom=16mm, left=14mm, right=14mm}

\examsetup{
  question/show-answer = true,
  fillin/show-answer = true,
  solution/show-solution = show-stay,
}

% ---------- 教师版解答样式 ----------
\newtcolorbox{solutionblock}{
  enhanced, breakable,
  colback=edu-blue-bg,
  colframe=edu-blue!25,
  boxrule=0pt,
  borderline west={2pt}{0pt}{edu-blue!50},
  arc=0pt,
  left=3mm, right=3mm, top=2mm, bottom=2mm,
  before skip=2mm,
  after skip=3mm,
  fontupper=\small,
}

% 解答步骤编号
\newcounter{solstep}
\newcommand{\solstep}[2]{%
  \stepcounter{solstep}%
  \par\medskip
  \noindent\hangindent=6mm
  {\color{edu-blue}\bfseries\textcircled{\scriptsize\thesolstep}\enspace #1}\enspace #2\par
}

\begin{document}

\section*{通过交点、面积、距离等条件求点坐标（三） · 练习卷（教师版）}

\section*{一、选择题（每题 12 分, 共 48 分）}


\begin{keyideabox}
  \textbf{中等思想：绝对值与高线转换}
  当三角形的底平行于坐标轴时，面积 $= \frac{1}{2} \times \text{底} \times |\text{高}|$。高线等于另一坐标的分量之差的绝对值，会产生正负两解。
\end{keyideabox}



\begin{question}[points=12]
  已知点 $O(0, 0)$，点 $A(4, 0)$。若点 $P$ 在直线 $y = x$ 上，且 $\triangle OAP$ 的面积为 $6$，则点 $P$ 的纵坐标为
  \begin{choices}
    \item 3
    \item -3
    \item 3 或 -3
    \item 1.5 或 -1.5
  \end{choices}
  \paren[C]
  \begin{solutionblock}
    底 $OA = 4$ 在 $x$ 轴上，高为 $|y_p|$。$S = \frac{1}{2} \times 4 \times |y_p| = 2|y_p| = 6 \implies |y_p| = 3 \implies y_p = \pm 3$。
  \end{solutionblock}
\end{question}



\begin{question}[points=12]
  在直线 $y = -2x + 6$ 与坐标轴围成的 $\triangle AOB$ 中，点 $C$ 是 $OA$ （$x$ 轴上交点）上的一点。若 $\triangle OBC$ 的面积为 $\triangle AOB$ 的 $\dfrac{1}{3}$，则 $OC$ 的长为
  \begin{choices}
    \item 1
    \item 2
    \item 3
    \item 1.5
  \end{choices}
  \paren[A]
  \begin{solutionblock}
    $A(3, 0)$，$B(0, 6)$。$OA = 3$。$\triangle OBC$ 与 $\triangle AOB$ 共用高 $OB$。面积比即为底之比，即 $OC = \frac{1}{3} OA = 1$。
  \end{solutionblock}
\end{question}



\begin{question}[points=12]
  点 $P$ 在直线 $y = 2x - 4$ 上。若 $\triangle OPA$ 的顶点为 $O(0, 0)$，$A(3, 0)$，且其面积为 $3$，则点 $P$ 的纵坐标绝对值为
  \begin{choices}
    \item 2
    \item 4
    \item 3
    \item 1
  \end{choices}
  \paren[A]
  \begin{solutionblock}
    底 $OA = 3$，面积 $= \frac{1}{2} \times 3 \times |y_p| = 3 \implies |y_p| = 2$。
  \end{solutionblock}
\end{question}



\begin{question}[points=12]
  已知 $A(0, 4)$，$B(6, 0)$。点 $C$ 在第一象限内，若 $AC \perp y$ 轴且 $AC = 3$，则点 $C$ 的坐标为
  \begin{choices}
    \item $(3, 4)$
    \item $(4, 3)$
    \item $(3, 0)$
    \item $(0, 3)$
  \end{choices}
  \paren[A]
  \begin{solutionblock}
    $AC \perp y$ 轴说明 $y_c = y_a = 4$；$AC = 3$ 且在第一象限，故 $x_c = 3$。坐标为 $(3, 4)$。
  \end{solutionblock}
\end{question}


\section*{二、填空题（每题 12 分，共 12 分）}


\begin{question}[points=12]
  点 $P$ 在直线 $y = 2x - 4$ 上，$\triangle OPA$ 中 $O(0, 0)$，$A(3, 0)$。若其面积为 $5$，且点 $P$ 在第一象限内，则点 $P$ 的坐标为 \fillin。
  \paren[$(\dfrac{11}{3}, \dfrac{10}{3})$]
  \begin{solutionblock}
    底 $OA = 3$，高为 $|y_p|$。面积 $= \frac{1}{2} \times 3 \times |y_p| = 5 \implies |y_p| = \frac{10}{3}$。在第一象限，故 $y_p = \frac{10}{3}$。代入直线方程：$\frac{10}{3} = 2x - 4 \implies 2x = \frac{22}{3} \implies x = \frac{11}{3}$。坐标为 $(\frac{11}{3}, \frac{10}{3})$。
  \end{solutionblock}
\end{question}


\section*{三、解答题（共 40 分）}

\needspace{8\baselineskip}

\begin{problem}[points=20]
  直线 $y = -\dfrac{4}{3}x + 4$ 与 $x$ 轴交于点 $A$，与 $y$ 轴交于点 $B$。点 $P$ 是直线 $AB$ 上的一个动点。

(1) 求点 $A$ 和点 $B$ 的坐标；
(2) 若点 $P$ 满足 $\triangle OAP$ 的面积为 $4$，求点 $P$ 的坐标。
  \answerarea[60mm]
  \setcounter{solstep}{0}
  \begin{solutionblock}
    \textbf{答案：}(1) $A(3, 0)$，$B(0, 4)$；(2) $P\left(5, -\dfrac{8}{3}\right)$ 或 $P\left(1, \dfrac{8}{3}\right)$\par\medskip
    \solstep{求 A、B 坐标}{令 $y = 0 \implies x = 3$，故 $A(3, 0)$；令 $x = 0 \implies y = 4$，故 $B(0, 4)$。}
    \solstep{利用面积建立方程}{底 $OA = 3$ 在 $x$ 轴上，高为点 $P$ 纵坐标的绝对值 $|y_p|$。\\ $S_{\triangle OAP} = \frac{1}{2} \times OA \times |y_p| = \frac{3}{2}|y_p| = 4 \implies |y_p| = \frac{8}{3}$。}
    \solstep{分情况求解}{第一种情况：$y_p = \frac{8}{3}$。代入 $y = -\frac{4}{3}x + 4$：\\ $\frac{8}{3} = -\frac{4}{3}x + 4 \implies -\frac{4}{3}x = -\frac{4}{3} \implies x = 1$。坐标为 $P_1\left(1, \frac{8}{3}\right)$。\\ 第二种情况：$y_p = -\frac{8}{3}$。代入直线方程：\\ $-\frac{8}{3} = -\frac{4}{3}x + 4 \implies -\frac{4}{3}x = -\frac{20}{3} \implies x = 5$。坐标为 $P_2\left(5, -\frac{8}{3}\right)$。}
  \end{solutionblock}
\end{problem}


\needspace{8\baselineskip}

\begin{problem}[points=20]
  在平面直角坐标系中，已知点 $A(0, 4)$，点 $C(5, 0)$，点 $B$ 在第一象限内，满足 $BA \perp y$ 轴。若 $AB = 6$。

(1) 求点 $B$ 的坐标及直线 $BC$ 的解析式；
(2) 若四边形 $ABCD$ 是等腰梯形，且 $AB \parallel DC$，$AD = BC$，求点 $D$ 的坐标。
  \answerarea[60mm]
  \setcounter{solstep}{0}
  \begin{solutionblock}
    \textbf{答案：}(1) $B(6, 4)$，$y = 4x - 20$；(2) $D(1, 0)$\par\medskip
    \solstep{求 B 坐标与 BC 解析式}{$BA \perp y$ 轴 $\implies y_b = y_a = 4$。$AB = 6$ $\implies x_b = 6$。故 $B(6, 4)$。\\ 设 $BC: y = kx + b$，代入 $C(5, 0)$ 和 $B(6, 4)$：\\ $\begin{cases} 5k + b = 0 \\ 6k + b = 4 \end{cases} \implies \begin{cases} k = 4 \\ b = -20 \end{cases}$，故 $BC: y = 4x - 20$。}
    \solstep{利用等腰梯形几何条件求 D}{因为 $AB \parallel DC$，$AB$ 水平，所以 $DC$ 也水平。$D$ 在 $x$ 轴上（因为 $C(5,0)$ 在 $x$ 轴上）。\\ 设 $D(x_d, 0)$。腰相等 $AD = BC$ $\implies AD^2 = BC^2$。\\ $BC^2 = (6-5)^2 + (4-0)^2 = 1 + 16 = 17$。\\ $AD^2 = (x_d - 0)^2 + (0 - 4)^2 = x_d^2 + 16$。\\ 令 $x_d^2 + 16 = 17 \implies x_d^2 = 1 \implies x_d = 1$ 或 $x_d = -1$。}
    \solstep{筛选解并写出结论}{当 $x_d = -1$ 时，四边形退化（不合梯形定义，且 $CD = 6 = AB$ 变成平行四边形）。\\ 故取 $x_d = 1$，点 $D$ 坐标为 $(1, 0)$。}
  \end{solutionblock}
\end{problem}


\clearpage
\section*{参考答案}


\begin{keyideabox}
  \textbf{选择题}
  1. C  2. A  3. A  4. A
\end{keyideabox}



\begin{keyideabox}
  \textbf{填空题}
  1. $\left(\dfrac{11}{3},\ \dfrac{10}{3}\right)$
\end{keyideabox}



\begin{keyideabox}
  \textbf{解答题}
  第 1 题：(1) $A(3, 0)$，$B(0, 4)$；(2) $P\left(5, -\dfrac{8}{3}\right)$ 或 $P\left(1, \dfrac{8}{3}\right)$
第 2 题：(1) $B(6, 4)$，$y = 4x - 20$；(2) $D(1, 0)$
\end{keyideabox}



\end{document}
